\chapter{The Socket-Style SDK and the Multiplexed Transfer}\label{ch:sdkmux}
\chaptermark{SDK and Multiplexed Transfer}

\section{The socket-style SDK}\label{sec:sdk}

\subsection{Contract}

\code{fuse/sdk.h} is a C API that will look familiar to anyone who has written a TCP server:
\code{fuse_listen}, \code{fuse_accept}, \code{fuse_connect}, \code{fuse_send}, \code{fuse_recv} and
\code{fuse_close}. Its semantics differ from TCP in four ways that a programmer must know:
\begin{itemize}
  \item \textbf{Messages, not a byte stream.} What one \code{fuse_send} delivers, exactly one \code{fuse_recv}
    returns. Boundaries are preserved, as with \code{SOCK_SEQPACKET}, so there is no length prefixing and no
    partial read.
  \item \textbf{Reliable and ordered over UDP.} Loss is repaired underneath; the application sees either the
    whole message or an error.
  \item \textbf{Completion means delivery.} \code{fuse_send} returns once the peer has acknowledged the message,
    not when it has been queued locally.
  \item \textbf{Optional encryption.} The same pre-shared key on both ends seals every message with AES-256-GCM.
\end{itemize}
A message may be up to 64~MiB. A connection may be used from several threads, typically one sending and one
receiving, and serialises its own operations internally.

\subsection{Connecting: one listen port, a fresh port per connection}

Figure~\ref{fig:sdkconnect} shows how a connection is established:
\begin{enumerate}
  \item The client sends HELLO($\mathit{salt}_c$, challenge) to the server's well-known port.
  \item The server allocates a connection object, a fresh ephemeral socket and a pump thread.
  \item It replies with HELLO\_ACK($\mathit{salt}_s$, proof) \emph{from the new socket}, and the client adopts
    that source address as its peer.
\end{enumerate}
As in TFTP~\cite{rfc1350}, a single open port serves any number of connections without connection identifiers.

The session key is derived with HKDF from the pre-shared key and the concatenation
$\mathit{salt}_c\,\|\,\mathit{salt}_s$, so that neither side alone determines it. The server's salt was added
in version~4 after a subtle problem was found. With only the client's salt, a HELLO retried before its first
reply arrived could create two live connections with the same key, and both would start their sequence numbers
at zero: an AES-GCM (key, nonce) collision. The proof is the client's challenge sealed under the session key, and
it demonstrates that the server holds the key. The two directions of a connection use different stream
identifiers (0 for client to server, 1 for server to client), and because the stream identifier is part of the
nonce, the directions never share one.

\begin{figure}[H]
\centering
\fitwidth{\begin{tikzpicture}[font=\small,yscale=0.85]
  \node[box] at (0,0.4) {client (ephemeral port)};
  \node[box] at (6,0.4) {server listen port};
  \node[box] at (11,0.4) {server per-connection port};
  \foreach \x in {0,6,11} \draw[thick,black!40] (\x,0) -- (\x,-5.2);
  \draw[msg] (0,-0.7) -- node[above,sloped]{HELLO($\mathit{salt}_c$, challenge)} (6,-1.2);
  \node[lbl,right] at (6.1,-1.6) {allocate connection,\\socket and pump thread};
  \draw[msg] (11,-2.4) -- node[above,sloped,pos=0.6]{HELLO\_ACK($\mathit{salt}_s$, proof)} (0,-3.1);
  \draw[msg,<->] (0,-3.9) -- node[above]{DATA / ACK / NACK} (11,-3.9);
  \draw[msg,<->] (0,-4.8) -- node[above]{CLOSE} (11,-4.8);
\end{tikzpicture}}
\caption{SDK connection establishment. The client adopts the HELLO\_ACK's source address as its peer.}\label{fig:sdkconnect}
\end{figure}

\subsection{Inside a connection}

Each connection has one UDP socket, one peer and a background \emph{pump} thread that does all the wire work:
receiving, reassembling, acknowledging and retransmitting. Because of the pump, a blocked \code{fuse_recv} still
sends acknowledgements and a blocked \code{fuse_send} still repairs loss. The application calls are thin. They
hand a message to the pump, or take one from it, through a mutex and a condition variable.

\paragraph{Framing.} A message is carried as a run of DATA blocks whose \code{offset} is the byte position
within the message, with LastBlock on the final one. The receiver knows the message is complete when its window
base passes the final block, so a truncated message can never reach the application. Block $\mathit{seq}$ of a
message must sit exactly at offset $(\mathit{seq}-\mathit{first\_seq})\times1200$, so a forged block cannot claim
an offset near 64~MiB and force a large allocation (R8).

\paragraph{Backpressure.} Completed messages wait in an inbox capped at 256~MiB of undelivered data. When the cap
is reached, new blocks are not admitted, and the sender's own retransmission logic stalls until the application
drains the inbox.

\paragraph{Liveness\vviii.} A connection that has sent nothing for one second sends a HEARTBEAT. A peer that
has sent nothing at all (no data and no authentic control message) for $\max(\mathit{timeout\_ms},3\,\mathrm{s})$
is declared lost, and blocked calls return \code{FUSE_ERR_TIMEOUT}.

\paragraph{Closing.} \code{fuse_close} sends CLOSE and stops the pump. The connection is reference-counted by
the application threads inside \code{fuse_send} and \code{fuse_recv}, and the last one to leave frees it, so
closing can never free a connection another thread is still using.

\subsection{Security}

The SDK was left out of the first version~8 pass and brought up to the same standard by the second. It now
shares the bulk transfer's protections:
\begin{itemize}
  \item retransmissions are sealed with the retransmission bit outside the AAD, which removes the nonce-reuse
    flaw (Section~\ref{sec:nonce});
  \item data keys rotate every $2^{22}$ blocks, and blocks outside the receive window are dropped before
    decryption, so forged far-off sequence numbers cannot force key derivations (R7);
  \item under a pre-shared key, ACK, NACK, CLOSE, HEARTBEAT and the path probes are signed and replay-checked;
    in plaintext they are acted on only from the peer's current address, where previously a two-byte CLOSE from
    anywhere ended a connection;
  \item the peer address moves only after the new source answers a path challenge, using the shared
    \code{PathValidator};
  \item binary keys with zero bytes are supported through \code{pre_shared_key_len}.
\end{itemize}

\noindent Three issues remain open, listed in Table~\ref{tab:open}:
\begin{itemize}
  \item the server still allocates a connection, a socket and a thread for every HELLO before the client
    proves anything; heartbeats now bound how long such a connection lives;
  \item the SDK is stop-and-wait per message;
  \item its registry uses about 17~MB per connection, because each slot embeds a 16~KiB payload while the SDK
    only sends 1,200-byte blocks.
\end{itemize}

\subsection{HTTP over the SDK}

\code{fuse/http.h} maps one HTTP/1.1 request and one response onto one SDK message each way, with discovery
through the Alt-Svc token \code{hfuse-00}~\cite{rfc7838}. Since the second pass it is validated and bounded:
\begin{itemize}
  \item methods, paths, headers and reason phrases are checked against the HTTP grammar, so CR/LF injection is
    refused;
  \item \code{timeout_ms} is a budget for the whole exchange;
  \item the server answers a malformed request with 400 and an invalid handler response with 500, gives each
    connection \code{timeout_ms} to send its request, serves at most 256 connections at once, and waits for
    in-flight handlers before returning;
  \item a dedicated test suite, \code{test_http}, covers this behaviour.
\end{itemize}

\section{The multiplexed transfer}

\subsection{Purpose}

The bulk transfer ties one thread and one port to each lane. The multiplexed transfer
(\code{fuse/mux_transfer.hpp}) instead carries $N$ logical streams over \emph{one} socket, serviced by a pool of
worker threads whose size adapts to the load (Section~\ref{sec:orch}). It is the right choice when an
application wants more logical streams than operating-system threads.

\subsection{Session timeline}

The transfer proceeds in five steps:
\begin{enumerate}
  \item \textbf{SETUP.} The sender opens one socket and runs the SETUP handshake (Section~\ref{sec:setup}),
    offering $N$ streams. Every stream is \code{LOSSLESS}, with the configured block and window sizes and stream
    identifiers $0,\ldots,N-1$.
  \item \textbf{Validation.} The receiver validates the offer and, in particular, requires the identifiers to be
    exactly $0,\ldots,N-1$\vviii, because it indexes its stream table with them.
  \item \textbf{Streams.} Both sides build a fixed table of per-stream state. Each stream then runs its own
    StreamStart, DATA, ACK and NACK exchange, with the shard split as in the bulk transfer.
  \item \textbf{Workers.} The workers service the streams until all are done. A control thread calls
    \code{tick} on the orchestrator and polls for completion and timeouts.
  \item \textbf{Errors.} A version-mismatch reply during SETUP is reported as \code{VersionMismatch}\vviii{}
    rather than a timeout.
\end{enumerate}

\subsection{What a worker does}

Every worker runs the same task. It first tries, without blocking, to take the socket's receive lock. If it gets
the lock, it receives a batch of datagrams and dispatches each one to its stream:
\begin{itemize}
  \item an ACK with the right nonce is applied with \code{apply_ack}\vviii;
  \item a NACK queues the missing blocks on the stream's retransmission backlog, de-duplicated\vviii;
  \item a duplicate SETUP reply is answered with another FINACK.
\end{itemize}
It then visits the streams round-robin, starting from a shared cursor, and performs one unit of work on the first
stream that has any. In order of priority:
\begin{enumerate}
  \item (re)send the StreamStart, every 50~ms until acknowledged;
  \item send one queued retransmission;
  \item handle an expired RTO by resending the base and known holes\vviii;
  \item send one new block if the window allows.
\end{enumerate}
Each stream has its own mutex, so any worker may service any stream. The table of streams is never resized
after SETUP, so indexing it needs no lock.

\subsection{Per-stream state machine}

A sending stream moves through three states:
\begin{itemize}
  \item \emph{Starting}: StreamStart is sent every 50~ms until any ACK carrying the stream's nonce arrives.
  \item \emph{Sending}: the window is filled, and NACKs and timeouts are serviced.
  \item \emph{Done}: the base has reached the next sequence number and the last block has been sent.
\end{itemize}
The receiving side mirrors the bulk-transfer receiver without the pipeline. Blocks are copied into a per-stream
buffer under the stream's mutex, and the stream is done when its base passes the last block.

\subsection{Differences from the bulk transfer}

The multiplexed transfer lacks several features of the bulk transfer:
\begin{itemize}
  \item encryption, and with it control-plane authentication;
  \item address validation and NAT re-anchoring;
  \item adaptive block size;
  \item segmentation offload;
  \item resume.
\end{itemize}
It shares the version~8 reliability improvements: selective acknowledgement, the RFC~6298 timer with backoff,
slow start (starting at 16 blocks), and send times with a random offset. Several version~8 robustness fixes were
specific to it:
\begin{itemize}
  \item a NACK for a block already waiting in the backlog is not queued twice;
  \item a full backlog causes an immediate resend instead of a dropped request;
  \item the stream identifiers are validated.
\end{itemize}
